\section{Ising Model in Two Dimensions}

We will not solve exactly the 2D Ising model, but we will give some arguments about its behavior and the nature of the phase transition it exhibits, and then in the end we will present the proof of the existence of a phase transition in the 2D Ising model at finite temperature, which is a non-trivial result that highlights the importance of dimensionality and fluctuations in determining the behavior of the system near critical points.

\subsection{Phase Transition - Euristic}

Consider the 2D Ising model on a square lattice, where each spin interacts with its four nearest neighbors. Take the limit for a very low temperature \(T \to 0\), and suppose that there is a region in the first 1D raw where the system is in a fully magnetized state, with all \(l\) spins aligned (e.g., \(m = +1\)). We want to ask ourselves if this state is stable against the formation of a droplet of flipped spins (e.g., \(m = -1\)) of size \(l \times l\). The energy cost of creating such a droplet can be estimated as follows:
\begin{itemize}
    \item Case A: \(\uparrow\,\uparrow\,\dots \,\uparrow \), where we have a single domain wall of length \(l\) separating the two regions of opposite magnetization. The energy cost of creating this domain wall is proportional to its length, which gives
          \[
              E = - J l,
          \]
          where \(J\) is the coupling constant. The energy cost is negative, which means that it is energetically favorable to create such a droplet, and the system will tend to form domains of flipped spins, leading to a disordered state with zero magnetization.
    \item Case B: \(\uparrow\,\dots \,\uparrow \,\downarrow\,\dots \,\downarrow \), where we have two domain walls of length \(l\) separating the three regions of magnetization. The energy cost of creating these two domain walls is proportional to their total length, which gives
          \[
              E = - J l + 2 J l = J l,
          \]
          where the first term corresponds to the energy cost of creating the first domain wall, and the second term corresponds to the energy cost of creating the second domain wall. The energy cost is positive, which means that it is energetically unfavorable to create such a droplet.
\end{itemize}

\begin{remark}[entropy of domains]
    ...
\end{remark}

If we now compute the free energy cost of creating such a droplet, we need to take into account both the energy cost and the entropy gain associated with the formation of the droplet. The free energy cost can be expressed as
\[
    F_B - F_A = \Delta E - T \Delta S = 2 J l - T k_B \log(l-1),
\]
since the energy cost of creating the droplet in case B is \(2 J l\) and the entropy gain is proportional to the logarithm of the number of ways to arrange the droplet, which is approximately \(\log(l-1)\). As \(l\) increases, the energy cost dominates over the entropy gain, and the free energy cost becomes positive, which means that it is energetically unfavorable to create such a droplet. This indicates that the system will tend to remain in a magnetized state with non-zero magnetization, and it suggests that there is a phase transition at a finite temperature \(T_c\) below which the system exhibits spontaneous magnetization.

This argument provides an intuitive understanding of why the 2D Ising model exhibits a phase transition at finite temperature, and it highlights the importance of dimensionality and fluctuations in determining the behavior of the system near critical points.

It seems that breaking a domain wall is energetically unfavorable, but there is a maximum value of \(l\) for which the entropy does not make unfavorable the creation of the droplet, and this maximum value of \(l\) can be estimated by setting the free energy cost to zero:
\[
    \Delta F = 0 \implies 2 J l - T k_B \log(l-1) = 0 \implies l \sim e^{\frac{2 J}{T k_B}} = \xi,
\]
which indicates that there is a characteristic length scale \(l\) above which the domain wall are not favorable anymore, and this length scale can be identified with the correlation length \(\xi\) of the system. As the temperature approaches the critical temperature \(T_c\) from below, the correlation length diverges, indicating that the system exhibits long-range correlations and critical behavior at the phase transition.

Euristic arguments in statistical physics can provide valuable insights into the behavior of complex systems and help to understand the underlying mechanisms driving phase transitions and critical phenomena. Sometimes, when a problem is mathematically intractable, we can use physical intuition and qualitative reasoning to make predictions about the behavior of the system, which can then be used to address the problem in new ways and in the end find the analutical solution.

\subsection{Peierls's argument - Euristic}

Let us compare the free energy of the two configurations:\TODO{add picture of the two configurations}
\[
    \dots
\]
We ask what is the difference in energy of the two 2D configurations described. Each set of nearest neighbours of the subdomain considered in the configuration contributes to the change in energy from one configuration to the other. In particular
\[
    \Delta U = U_B - U_A = 2 J l,
\]
where now in 2D \(l\) represents the lenght of the perimeter of the flipped subdomain, which is proportional to the number of nearest neighbor pairs that are affected by the change in configuration. The energy cost of creating a droplet of flipped spins is proportional to the length of the domain wall, which is \(2 J l\) in this case. This energy cost is positive, which means that it is energetically unfavorable to create such a droplet, and it suggests that the system will tend to remain in a magnetized state with non-zero magnetization, indicating the presence of a phase transition at finite temperature.

But we have always to take into account the entropy gain associated with the formation of the droplet, which can be estimated after counting the number of ways to arrange the droplet of flipped spins. The number of ways to arrange a droplet of size \(l\) can be approximated as the path we can take to create the domain wall, which is given by three choices of direction at each step (since we cannot go back),\footnote{Extimating like this is exagerated, since we are counting also ``snakes'' that do not close on a region, but let us do this anyay and then reason on it.} thus we can estimate the number of configurations as
\[
    \# DM \sim 3^l,
\]
thus if the number of domains wall of length \(l\) is given by \(\# DM\), then the entropy gain associated with the formation of the droplet can be estimated as
\[
    \Delta S = k_B \log(\# DM) = k_B l \log(3),
\]
which is proportional to the length of the domain wall.

Thus in the end, the free energy cost of creating such a droplet can then be expressed as
\[
    \Delta F = \Delta U - T \Delta S = 2 J l - T k_B l \log(3) = l (2 J - T k_B \log(3)),
\]
which indicates that there is a critical temperature \(T_c\) at finite temperature, at which we observe a phase transition.

\subsection{Rigorous Proof of Phase Transition}\TODO{Add reference to 14.3 Huang's book stat mech.}

Let us consider a system with \textbf{fixed open boundary conditions}, where we have an Ising model defined on a square lattice of size \(N \times N\), and we fix the spins on the boundary to be all up (i.e., \(s_i = +1\) for all spins on the boundary). We want to show that there is a phase transition at finite temperature, which can be done by comparing the free energy of two configurations: one with all spins up (the ordered phase) and one with a droplet of flipped spins (the disordered phase).

The hamiltonian of the system is unchanged, and it is given by
\[
    H = -J \sum_{\langle i,j \rangle} s_i s_j - h \sum_{i} s_i,
\]
where the first sum is taken over nearest neighbor pairs of spins, and the second sum is taken over all spins in the lattice.

We can fix the temperature and compute the partition function and magnetization for both configurations. Then in the end we will find: if we denote by \(N_-\) the average number of down spins, we will have only two cases:
\begin{itemize}
    \item \(\lim_{N \to \infty} \frac{N_-}{N} = \frac{1}{2}\): in this case the system is in a disordered phase with zero magnetization, since the number of up spins and down spins are equal in the thermodynamic limit.
    \item \(\lim_{N \to \infty} \frac{N_-}{N} < \frac{1}{2}\): in this case the system is in an ordered phase with non-zero magnetization, since the number of up spins is greater than the number of down spins in the thermodynamic limit.
\end{itemize}
We used a preferred direction setting the boundary conditions, but we could have done the same analysis with all spins down at the boundary, and we would have found the same result. Exact numbers will fluctuate, but these ratios will converge to the same limit, and in the TD limit we will have a well defined phase transition at finite temperature.